% ---------------------------------------------------------------------------
% OPTIONAL: Ergänzung für Kapitel 8 (Validierung und Ergebnisse).
% Abschnitt 8.1 sagt "Die Starship-Ergebnisse sind separat zu ergänzen". Dieser Abschnitt
% wird als NEUER Abschnitt 8.5 am Ende von Kapitel 8 angehängt; die Clevon-Abschnitte
% 8.1 bis 8.4 bleiben unverändert.
% ---------------------------------------------------------------------------
\section{Starship-Modell}
\label{sec:val_starship}

Die Ergebnisse des Starship beruhen auf skriptgesteuerten Fahr- und Bordsteintests in
Town10HD, dem Auslesen der Physikwerte im laufenden Spiel und den ROS-2-Gesamttests
aus Kapitel~7. Rohdaten und Prüfskripte liegen im Projekt-Repository unter
\texttt{02\_starship\_emodell/docs/messungen} und \texttt{02\_starship\_emodell/scripts/carla}
\parencite{rauch_lee_repo}.

Der Abnahmetest liest zunächst die Blueprint-Attribute und die Physikwerte des
gespawnten Fahrzeugs aus und vergleicht sie mit den Sollwerten. Bestätigt wurden
35\,kg, sechs Räder mit 9\,cm Radius, Lenkwinkel nur an der Vorderachse, Antrieb aller
sechs Räder, 720\,min$^{-1}$ und 3\,Nm. Anschließend fährt der Test definierte Manöver
(Tabelle~\ref{tab:val_starship}).

\begin{table}[htbp]
  \centering
  \caption{Abnahmetest des Starship-Modells (Town10HD, CARLA 0.10.0); eigene Messung vom
  20.10.2026.}
  \label{tab:val_starship}
  \begin{tabular}{ll}
    \hline
    Prüfung & Ergebnis \\
    \hline
    Ruhelage & steht auf allen sechs Rädern \\
    Höchstgeschwindigkeit, Vollgas 8\,s & 6,27 bis 6,31\,km/h, 90\,\% nach 0,8\,s \\
    Nick-/Rollwinkel geradeaus & unter 3° \\
    Bremsweg aus ca. 6\,km/h & 0,22\,m \\
    Volllenkung bei 60\,\% Gas & Wenderadius 2,2 bis 2,4\,m, Rollwinkel unter 3° \\
    Rückwärtsfahrt & 6,1\,km/h \\
    Handbremse & bremst bis zum Stillstand \\
    Bordstein 15\,cm hoch / runter & 10 von 12 / 8 von 8 (Raycast, Tabelle~\ref{tab:starship_bordstein}) \\
    \hline
  \end{tabular}
\end{table}

Alle 26 Sensoren lieferten im ROS-2-Gesamttest gleichzeitig Daten mit ihrer Sollrate;
nur die IMU war durch den Simulationstakt begrenzt (Tabelle~7.3). Ultraschallsensoren
vorne erkannten ein Testfahrzeug in 1,15 bis 2,27\,m Abstand. Die Steuerung über ROS~2
erreichte die Höchstgeschwindigkeit, fuhr eine Linkskurve und bremste bis zum
Stillstand (Tabelle~7.2).

Die Prüfungen belegen die Funktionsfähigkeit des Modells, keine Übereinstimmung mit dem
realen Fahrzeug. Messwerte des Originals zu Beschleunigung, Bremsweg, Wendekreis und
Kletterverhalten lagen nicht vor. Das Herunterfahren vom Bordstein erzeugt Stoßspitzen
von etwa 4\,g, und etwa jede sechste Bordsteinauffahrt bleibt hängen. Die Bordsteinmessung mit der aktuellen Spherecast-Radabfrage
sowie ein Test mit synchronem Simulationstakt, Steigungen und unterschiedlichen Bordsteinhöhen stehen aus.
